\documentclass[11pt]{article}
\usepackage[margin=1in]{geometry}
\usepackage{amsmath,amssymb,booktabs,enumitem,hyperref}
\hypersetup{colorlinks=true,urlcolor=blue,linkcolor=blue}
\setlength{\parindent}{0pt}
\setlength{\parskip}{0.6em}
\setlist{nosep,leftmargin=*}
\title{Learning to collect evidence against adversarial tests}
\author{Experiment explainer and protocol}
\date{17 November 2026}
\begin{document}
\maketitle

\textbf{Question.} Can a learned generator find interactions that expose missing
information, and thereby help a collector acquire more reusable evidence than
adding random target interactions? This pilot tests that idea on small environments
where an independent calculation can distinguish learning failure from an
unavoidable information tradeoff.

\section{The algorithm should accommodate tradeoffs}
Suppose an environment has two hidden bits but allows reading only one of them.
A collector that reads the left bit can reproduce another left-bit reading
exactly. A right-bit target exposes information that collector lacks. Reading
right instead loses the previous capability; no collector can reproduce both
targets exactly. Randomly choosing which bit to read balances the errors, but
still reveals only one bit on each run.

We will therefore minimize the \emph{largest simulation error} across retained
targets. Each error is measured using the best available randomized decoder,
required to work across the possible worlds. The algorithm may trade a smaller
error on one target for a larger error on another. Positive error can be the best
feasible result. Earlier targets remain in the archive so that their losses
remain visible and continue to enter the objective.

\section{Why maximum and the reference score differ}
The paper's reference score $J_w$ subtracts a positively weighted average of
eventual target errors from one. A decrease in any error, with none increasing,
strictly improves that weighted score. A largest-error score can stay unchanged
when a smaller error decreases. Its preference for balancing the hardest targets
is a reasonable additional choice, but does not retain every strict improvement.

For a concrete native example, let three hidden bits vary independently across
eight worlds. One irreversible action reveals the first bit, another reveals
the first two bits, and a third reveals the last bit. No later observation gives
information. Reading two bits strictly improves on reading the first bit alone,
yet both collectors have worst native simulation error $1/2$: neither reveals
the last bit. Mixing the three available target actions cannot increase that
bound. Thus taking the supremum over all native targets still allows this tie.

A finite maximum is continuous as a function of its finite vector of errors,
although ties can make it nondifferentiable. An infinite supremum need not be
continuous under coordinatewise convergence. The eventual collection-time limit
introduces another issue. This pilot uses finite budgets and finitely many
targets; it does not approximate the numerical value of $J_w$ with a proved error
bound.

\section{The small environments and available data}
Each world specifies two bits, giving four possible worlds. There are three
actions: read left, read right, and request a fresh fair coin. Every observation
is binary. A collector receives two action--observation pairs on each evaluated
run. Target interactions also last two steps.

We will use three environments:
\begin{enumerate}
\item \textbf{Compatible readings.} Reading an indicated bit is exact and can be
repeated. The collector can read both bits within the budget. The coin action
has no effect on later readings.
\item \textbf{Irreversible choice.} The first action returns the chosen bit or
a fair coin. Every second observation is zero. The two reading targets compete;
their best balanced error is $1/4$.
\item \textbf{Noisy readings.} Each bit reading is independently correct with
probability $0.85$. The coin action remains independent of the world. At this
finite budget, information can remain missing even after the best collection.
\end{enumerate}

The training procedure receives reset episodes from each of the four labelled
worlds, with uniformly random actions. From those samples it estimates the
observation probabilities after each history and proposed action. Training then
uses these empirical response tables. The trainer knows which world generated
each episode; the collector and decoder receive no world label. Each evaluated
record contains only the two interactions of the frozen collector. Training
episodes are not silently added to that record.

This is a sampled-data, model-based pilot with strong access assumptions:
repeatable resets and labelled training access to every world. It does not model
learning from one irreversible encounter with an unknown environment. The true
response probabilities and exact optimization routines are reserved for
evaluation, not used for learner updates or target selection.

\section{The three learned components}
A deterministic two-step target chooses an initial action and a second action
for each of the two possible first observations. There are $3^3=27$ such target
policies. The generator learns three categorical action distributions and uses
them to propose these contingent policies. This finite parameterization is a
deliberate first test; it supplies no generalization to longer programs.

The collector is a stochastic table over the seven decision histories: the
empty history and the six possible first action--observation pairs. The decoder
is a stochastic table whose input is one of the 36 possible full collected
histories. For each target, the decoder outputs one of its four observation
strings. The known deterministic target reconstructs its own actions from those
observations, so this output represents the full target history exactly.

Let $E_\pi(\theta,h)$ be the empirical probability of collected history $h$
under collector $\pi$ in world $\theta$. For a target $\tau$, let
$F_\tau(\theta,y)$ be the empirical probability of its observation string $y$.
The decoder $G_\tau(y\mid h)$ gives the probability of simulated output $y$
after collected history $h$. Its achieved error is
\[
 d(\pi,G_\tau,\tau)=\max_\theta\frac12\sum_y
 \left|\sum_h E_\pi(\theta,h)G_\tau(y\mid h)-F_\tau(\theta,y)\right|.
\]
Minimizing this error over all stochastic decoders gives finite-experiment
deficiency. A fitted decoder can have higher error because training is
imperfect. We will report that difference separately.

The collector and its decoder bank train to reduce the largest achieved error
over the retained targets. Updates alternate between fitting
decoders for the current collector and improving the collector with the
decoders held fixed. Categorical probabilities are represented by trainable
parameters; gradients of the finite empirical distribution loss provide the
updates. This optimization may find a local solution or fail to settle.
Each collector action retains probability at least $0.005$. We report both the
unrestricted reference optimum and the optimum with this same probability
floor, so the imposed randomization is not mistaken for failed optimization.

\section{One iteration and its comparison}
Both methods start with the repeated-left target and the same collector
initialization. At each iteration:
\begin{enumerate}
\item Fit decoders and update the collector against every archived target.
\item Freeze the collector and propose a bounded batch of omitted targets.
Fit their decoders with a separate fixed training budget.
\item The adversarial generator uses fitted simulation errors to update its
action probabilities and adds a difficult candidate. The random control samples
uniformly and adds a random candidate. Both receive the same number of candidate
queries and decoder updates; generator overhead and wall time are reported.
\item Refit the collector against the enlarged archive. Save the collector,
decoders, generator, losses and random seed, then perform the independent audit.
\end{enumerate}

A fair-coin target should be simulable by generating a fresh fair coin. Comparing
distributions makes that possible; predicting the specific random outcome of a
separate run would reward the wrong behavior. An undertrained decoder can still
make noise look difficult, so decoder fitting is part of the experiment rather
than an assumption that all observed error is missing information.

An additional control samples candidates uniformly but retains the candidate
with largest fitted error. This separates learning the proposal distribution
from selecting a difficult member of a random batch. A further diagnostic
trains against all 27 targets from the start. That diagnostic receives the same
number of collector updates but more decoder work, so it is not an equal-cost
competitor.

\section{Evaluation and the meaning of success}
The independent audit computes true-world deficiencies to all 27 targets using
linear programming. A separate optimization supplies the best finite minimax
collector, and numerical feasibility and dual bounds check the reference. These
calculations do not choose targets or train the learned collector.

We will compare adversarial and random target additions using paired seeds,
identical training samples, collection budgets, candidate budgets and collector
update counts. The primary measure is the true largest target deficiency after
each archive expansion. We also retain every target error, achieved decoder
error, distance above the reference optimum, action probabilities, selected
targets and computational cost. The main comparison includes every registered
seed and environment, including failures. Improvements on individual targets
will not be described as universal dominance.

Mechanics checks precede the main runs: normalized experiments; correct action
reconstruction; identical decoder and full-history evaluations; zero deficiency
for world-independent noise; the analytic irreversible bound; and agreement
between separate numerical evaluators. Development runs are labelled separately.
The final run configuration and source hashes will be frozen before the paired
main runs. Changes after inspecting main results require a separate recorded run.

The useful outcome could be better target discovery, an informative failure of
the learned loss, or evidence that the extra adversarial search is unnecessary
on these cases. We will not treat failure to find a new challenge as a convergence
certificate. This experiment supplies no universal stopping rule for the eventual
score, sample-efficiency theorem, or guarantee on unseen world classes.
\end{document}
